\chapter{Adding momentum: the LFISTA-style network}\label{ch:lfista}

\begin{careful}[Read the label]
This network is \textbf{our own combination}. It is inspired by LFISTA, but it is not the method of Moreau and
Bruna. Do not present it as an implementation of their paper. It is an exploratory run.
\end{careful}

\section{The motivation}

FISTA is the best method we have, so the question you asked was natural: can we merge the FISTA update with
learning? Is there literature? Yes, there is.

\section{LFISTA in the paper}

Moreau and Bruna, ``Understanding Trainable Sparse Coding via Matrix Factorization'', ICLR 2017, arXiv 1609.00285.
Their layer feeds in the last \emph{two} estimates, the way FISTA does:
\begin{equation}
z^{(k+1)} = \soft\big(W_g^{(k)} z^{(k)} + W_m^{(k)} z^{(k-1)} + W_e^{(k)} x\big).
\end{equation}
Here $x$ is the input signal, in their notation. Their $W$ matrices are dense and free. They start them at
the values that make $K$ layers equal $K$ steps of FISTA. And they train by minimising the \textbf{Lasso objective}
itself, not by comparing to a target.

\section{Why it cannot run as published}

Dense matrices between all code values at 0.1\um\ would need about 16 petabytes for each matrix. So you asked for
the closest feasible version, and that is what Job 3 is.

\section{What Job 3 does}

Take Eq.~20 and move the layer's starting point using momentum.
\begin{align}
y^{(k)} &= x^{(k)} + \beta_k\big(x^{(k)} - x^{(k-1)}\big),\\
x^{(k+1)} &= \soft_{\theta^{(k)}}\Big(y^{(k)} - \big(w^{(k)}\big)' * \big(\textstyle\sum_m d_m*y_m^{(k)} - b\big)\Big).
\end{align}

\begin{code*}{csc_lista.py}
            y = update if self.beta is None else update + self.beta[k] * (update - x)
            x = update
\end{code*}

Each $\beta_k$ is learned, starts at FISTA's own value, and stays within $[0,1]$. A test confirms that the untrained
network matches FISTA to within 2\% when filters are wide enough.

\begin{code*}{csc_lista.py}
            t, start_beta = 1.0, []
            for _ in range(layers):
                t_next = (1 + (1 + 4 * t * t) ** .5) / 2
                start_beta.append((t - 1) / t_next)
                t = t_next
\end{code*}

\subsection*{The link to LFISTA}

Expand the update and collect terms. Writing $W'$ for the correction filter and $D$ for the blurred dictionary:
\begin{equation}
x^{(k+1)} = \soft\Big(\underbrace{(1+\beta)(I-W'D)}_{W_g}\,x^{(k)}
+ \underbrace{(-\beta)(I-W'D)}_{W_m}\,x^{(k-1)}
+ \underbrace{W'}_{W_e}\,b\Big).
\end{equation}
So our layer has LFISTA's form, but with $W_g$, $W_m$ and $W_e$ all tied together through $W'$ and made convolutional.
That is far more restricted than their free matrices.

\subsection*{Training on the Lasso objective}

Like LFISTA, this loss never looks at the target. It only asks the codes to score well on the Lasso objective. The
best checkpoint, however, is still chosen by validation error against targets, the same as every other run.

\begin{code*}{csc_lista.py}
    def lasso_objective(self, b, lam):
        """Mean of 0.5 ||mask(Dz - b)||^2 + lam ||z||_1 per pixel, as in LFISTA training."""
        z = self.codes(b)
        operator, mask = self.physics(b.shape[-1])
        residual = mask * (operator.forward(z) - b)
        value = .5 * (residual ** 2).sum((1, 2)) + lam * z.abs().sum((1, 2, 3))
        return value.mean() / b.shape[-1] ** 2
\end{code*}

The value $\lambda=0.03$ is the one FISTA chose on validation. So in plain words, this network is being asked to
\emph{reach FISTA's answer in 32 steps}.

\section{How it differs from everything else}

\begin{center}
\begin{tabular}{lccc}
\toprule
 & Strict Eq.~20 & Job 3 & LFISTA as published \\
\midrule
Momentum & No & Yes, learned & Yes \\
Weights & Convolutional & Convolutional, tied & Dense, free \\
Loss & MSE with target & Lasso, no target & Lasso, no target \\
Problem size & $768\times768$ & $768\times768$ & Small \\
\bottomrule
\end{tabular}
\end{center}

\section{Result}

\begin{center}
\begin{tabular}{lccc}
\toprule
Method & Simulated validation & COMSOL & NRMSE \\
\midrule
FISTA-32, $\lambda=0.03$ (the start) & 50.9\% & 41.8\% & 3.00\% \\
\textbf{LFISTA-style, 32 layers} & \textbf{42.8\%} & \textbf{34.2\%} & \textbf{2.47\%} \\
FISTA-300, $\lambda=0.03$ (the goal) & 36.6\% & 32.1\% & 2.33\% \\
\bottomrule
\end{tabular}
\end{center}

Read this table as a story. The network starts as 32 steps of FISTA, and before training it scored 51.0\% on validation,
matching FISTA-32's 50.9\%. Training asked it to reach the same Lasso solution faster. It closed 57\% of the gap to
FISTA-300 on simulation and 78\% on COMSOL. Unlike the strict runs, it beat FISTA at equal depth on both.

\begin{center}
\begin{tabular}{rcc}
\toprule
Step & Validation error & Density \\
\midrule
0 & 51.0\% & 9.0\% \\
100 & 45.8\% & 18.4\% \\
500 & 43.3\% & 16.9\% \\
1,000 & 42.8\% & 16.6\% \\
\bottomrule
\end{tabular}
\end{center}

Its validation error improved all the way to the last step, so the best checkpoint is step 1,000. It has 24,601,664
parameters, which is strict 32 plus 32 values of $\beta$.

\begin{careful}[Three things to know]
\begin{itemize}
\item 15 of the 32 values of $\beta$ ended pressed against the upper limit of 1. That limit is our clamp, so training
may have wanted more momentum than we allowed.
\item It used $\lambda=0.03$, chosen before FISTA was retuned. Whether $\lambda=0.009$ helps has not been tried.
\item It is 0.4 points behind strict 64 layers with half the layers, but on one seed that gap means little. It remains our
own combination, so it is not the selected model.
\end{itemize}
\end{careful}
